\documentclass[dvipsnames]{article}

\usepackage{csquotes}

% Set page size and margins
% Replace `letterpaper' with`a4paper' for UK/EU standard size
\usepackage[letterpaper,top=2cm,bottom=2cm,left=3cm,right=3cm,marginparwidth=1.75cm]{geometry}

% Useful packages
\usepackage{amsmath}
\usepackage{graphicx}
\usepackage{minted2}
\usepackage{amsthm}
\usepackage{mathtools}
\usepackage[colorlinks=true, allcolors=blue]{hyperref}
\usepackage{fancyhdr}
\usepackage{enumerate}
\usepackage{tikz-cd}
\usepackage{mdframed}
\usepackage{multicol}

\setminted{frame=single}
\newmintinline[java]{java}{}

\usepackage{array}
\newcolumntype{L}[1]{>{\raggedright\let\newline\\\arraybackslash\hspace{0pt}}m{#1}}
\newcolumntype{C}[1]{>{\centering\let\newline\\\arraybackslash\hspace{0pt}}m{#1}}
\newcolumntype{R}[1]{>{\raggedleft\let\newline\\\arraybackslash\hspace{0pt}}m{#1}}

\DeclareMathOperator*{\argmin}{arg\,min}

\usepackage{tikz}
\usetikzlibrary{shapes}

\tikzset{
  treenode/.style = {align=center, inner sep=0pt, text centered,
    font=\sffamily},
  avl_node/.style = {treenode, circle, black, draw=black,
    text width=1.5em, very thick},% AVL node
  black_node/.style = {treenode, circle, white, font=\sffamily\bfseries, draw=black,
    fill=black, text width=1.5em},% red-black tree, black node
  red_node/.style = {treenode, circle, red, draw=red,
    text width=1.5em, very thick},% red-black tree, red node
  nil_leaf/.style = {treenode, rectangle, draw=black, fill=black,
    minimum width=0.3em, minimum height=0.3em}% red-black tree, nil
}

\newcommand{\haskell}[1]{\mintinline{haskell}{#1}}
\newcommand{\racket}[1]{\mintinline{racket}{#1}}
\newcommand{\prolog}[1]{\mintinline{prolog}{#1}}

\DeclareMathOperator{\xor}{\mathsf{xor}}
\DeclareMathOperator{\tru}{\mathsf{tru}}
\DeclareMathOperator{\fls}{\mathsf{fls}}
\DeclarePairedDelimiter{\ceil}{\lceil}{\rceil}

\begin{document}

\section*{Topic 7. Problem set (theory)}
\setcounter{section}{7}

\pagestyle{fancy}
\lhead{Solution authored by \textcolor{red}{Gabdulla Tukay \texttt{g.tukay@innopolis.university}}}
\rhead{}
\cfoot{Data Structures and Algorithms (Spring 2026), Innopolis University}

\theoremstyle{definition}
\newtheorem{problem}{Problem}[section]

\begin{problem}
  For each of the following equalities,
  either prove it using the axioms of probability and basic set identities~\cite[\S B.2--B.3]{CLRS-2022},
  \emph{or} provide a counterexample (a concrete sample space and events, with justification)
  to show that it does not hold in general.

  \begin{enumerate}
    \item \(\mathsf{Pr}\{A \setminus B\} = \mathsf{Pr}\{A\} - \mathsf{Pr}\{B\}\).

%%%%%%%%%%%%%%%%%%%%%%%%%%%%%%%%%%%%%%%%%%%%%%%%%
% Start of answer: Problem 1, Item 1
%%%%%%%%%%%%%%%%%%%%%%%%%%%%%%%%%%%%%%%%%%%%%%%%%
\color{purple}
\begin{mdframed}
  \color{purple}
  \emph{Replace with your proof or counterexample.}
\end{mdframed}
%%%%%%%%%%%%%%%%%%%%%%%%%%%%%%%%%%%%%%%%%%%%%%%%%
% End of answer: Problem 1, Item 1
%%%%%%%%%%%%%%%%%%%%%%%%%%%%%%%%%%%%%%%%%%%%%%%%%

\color{black}\noindent
    \item \(\mathsf{Pr}\{A \cup (B \cap C)\} = \mathsf{Pr}\{(A \cap B) \cup C\}\).

%%%%%%%%%%%%%%%%%%%%%%%%%%%%%%%%%%%%%%%%%%%%%%%%%
% Start of answer: Problem 1, Item 2
%%%%%%%%%%%%%%%%%%%%%%%%%%%%%%%%%%%%%%%%%%%%%%%%%
\color{purple}
\begin{mdframed}
  \color{purple}
  \emph{Replace with your proof or counterexample.}
\end{mdframed}
%%%%%%%%%%%%%%%%%%%%%%%%%%%%%%%%%%%%%%%%%%%%%%%%%
% End of answer: Problem 1, Item 2
%%%%%%%%%%%%%%%%%%%%%%%%%%%%%%%%%%%%%%%%%%%%%%%%%

  \end{enumerate}
\end{problem}

\newpage\color{black}
\begin{problem}
  For each of the following intuitive-looking equations, either prove it (by unwinding the definitions of random variables and expected value~\cite[\S B.2--B.3]{CLRS-2022}, and using properties such as linearity and independence where applicable), \emph{or} provide a counterexample (specifying the distribution of the involved random variables).

  Assume that all random variables take real values and that all expected values appearing in an equation exist.

  \begin{enumerate}
    \item \(E\!\left[\dfrac{X + Y + Z}{3}\right] = \dfrac{E[X] + E[Y] + E[Z]}{3}\).

%%%%%%%%%%%%%%%%%%%%%%%%%%%%%%%%%%%%%%%%%%%%%%%%%
% Start of answer: Problem 2, Item 1
%%%%%%%%%%%%%%%%%%%%%%%%%%%%%%%%%%%%%%%%%%%%%%%%%
\color{purple}
\begin{mdframed}
  \color{purple}
  \emph{Replace with your proof or counterexample.}
\end{mdframed}
%%%%%%%%%%%%%%%%%%%%%%%%%%%%%%%%%%%%%%%%%%%%%%%%%
% End of answer: Problem 2, Item 1
%%%%%%%%%%%%%%%%%%%%%%%%%%%%%%%%%%%%%%%%%%%%%%%%%

\color{black}\noindent
    \item \(E[\sqrt{XY}] = \sqrt{E[X]\,E[Y]}\) when \(X\) and \(Y\) are independent.

%%%%%%%%%%%%%%%%%%%%%%%%%%%%%%%%%%%%%%%%%%%%%%%%%
% Start of answer: Problem 2, Item 2
%%%%%%%%%%%%%%%%%%%%%%%%%%%%%%%%%%%%%%%%%%%%%%%%%
\color{purple}
\begin{mdframed}
  \color{purple}
  \emph{Replace with your proof or counterexample.}
\end{mdframed}
%%%%%%%%%%%%%%%%%%%%%%%%%%%%%%%%%%%%%%%%%%%%%%%%%
% End of answer: Problem 2, Item 2
%%%%%%%%%%%%%%%%%%%%%%%%%%%%%%%%%%%%%%%%%%%%%%%%%

\color{black}\noindent
    \item \(E[\max(X, Y)] = \max(E[X], E[Y])\).

%%%%%%%%%%%%%%%%%%%%%%%%%%%%%%%%%%%%%%%%%%%%%%%%%
% Start of answer: Problem 2, Item 3
%%%%%%%%%%%%%%%%%%%%%%%%%%%%%%%%%%%%%%%%%%%%%%%%%
\color{purple}
\begin{mdframed}
  \color{purple}
  \emph{Replace with your proof or counterexample.}
\end{mdframed}
%%%%%%%%%%%%%%%%%%%%%%%%%%%%%%%%%%%%%%%%%%%%%%%%%
% End of answer: Problem 2, Item 3
%%%%%%%%%%%%%%%%%%%%%%%%%%%%%%%%%%%%%%%%%%%%%%%%%

\color{black}\noindent
    \item \(E[X^3] = (E[X])^3\).

%%%%%%%%%%%%%%%%%%%%%%%%%%%%%%%%%%%%%%%%%%%%%%%%%
% Start of answer: Problem 2, Item 4
%%%%%%%%%%%%%%%%%%%%%%%%%%%%%%%%%%%%%%%%%%%%%%%%%
\color{purple}
\begin{mdframed}
  \color{purple}
  \emph{Replace with your proof or counterexample.}
\end{mdframed}
%%%%%%%%%%%%%%%%%%%%%%%%%%%%%%%%%%%%%%%%%%%%%%%%%
% End of answer: Problem 2, Item 4
%%%%%%%%%%%%%%%%%%%%%%%%%%%%%%%%%%%%%%%%%%%%%%%%%

  \end{enumerate}
\end{problem}

\newpage\color{black}
\begin{problem}[Modular-Search]
  Consider a linear-time search that traverses an array using \emph{modular arithmetic}: fix a constant step \(s\) (e.g.\ \(s = 17\)).
  Start at index \(1\), then repeatedly move by \(s\) modulo \(n\) (i.e., visit indices \(1,\; 1+s,\; 1+2s,\; \dots\), reducing modulo \(n\) so indices stay in \([1..n]\)).
  When \(\gcd(n,s) = 1\), every index is visited exactly once over \(n\) steps,
  so the algorithm is a linear scan in a fixed ``modular'' order.

  Now, assume the array \(A[1 \dots n]\) contains exactly \(k\) copies of the searched value \(x\) (and \(n - k\) other elements);
  the algorithm stops as soon as it finds an occurrence of \(x\).

  \begin{enumerate}
    \item Give a high-level description of this \textsc{Modular-Search} algorithm and write down clear pseudocode
    (including how the next index is computed using modular arithmetic).

%%%%%%%%%%%%%%%%%%%%%%%%%%%%%%%%%%%%%%%%%%%%%%%%%
% Start of answer: Problem 3, Item 1
%%%%%%%%%%%%%%%%%%%%%%%%%%%%%%%%%%%%%%%%%%%%%%%%%
\color{purple}
\begin{mdframed}
  \color{purple}
  \emph{Replace with your description and pseudocode.}
\end{mdframed}
%%%%%%%%%%%%%%%%%%%%%%%%%%%%%%%%%%%%%%%%%%%%%%%%%
% End of answer: Problem 3, Item 1
%%%%%%%%%%%%%%%%%%%%%%%%%%%%%%%%%%%%%%%%%%%%%%%%%

\color{black}\noindent
    \item Analyze the running time (number of comparisons until a copy of \(x\) is found, or \(n\) if none):
      \begin{enumerate}[(a)]
        \item[\textbf{worst-case}:] give the asymptotic running time in terms of \(n\) and \(k\), and briefly justify
        (e.g., how an adversary could arrange the \(k\) copies to maximize the running time).

%%%%%%%%%%%%%%%%%%%%%%%%%%%%%%%%%%%%%%%%%%%%%%%%%
% Start of answer: Problem 3, Item 2(a)
%%%%%%%%%%%%%%%%%%%%%%%%%%%%%%%%%%%%%%%%%%%%%%%%%
\color{purple}
\begin{mdframed}
  \color{purple}
  \emph{Replace with your worst-case analysis.}
\end{mdframed}
%%%%%%%%%%%%%%%%%%%%%%%%%%%%%%%%%%%%%%%%%%%%%%%%%
% End of answer: Problem 3, Item 2(a)
%%%%%%%%%%%%%%%%%%%%%%%%%%%%%%%%%%%%%%%%%%%%%%%%%

\color{black}\noindent
        \item[\textbf{best-case}:] give the asymptotic running time in terms of \(n\) and \(k\), and briefly justify.

%%%%%%%%%%%%%%%%%%%%%%%%%%%%%%%%%%%%%%%%%%%%%%%%%
% Start of answer: Problem 3, Item 2(b)
%%%%%%%%%%%%%%%%%%%%%%%%%%%%%%%%%%%%%%%%%%%%%%%%%
\color{purple}
\begin{mdframed}
  \color{purple}
  \emph{Replace with your best-case analysis.}
\end{mdframed}
%%%%%%%%%%%%%%%%%%%%%%%%%%%%%%%%%%%%%%%%%%%%%%%%%
% End of answer: Problem 3, Item 2(b)
%%%%%%%%%%%%%%%%%%%%%%%%%%%%%%%%%%%%%%%%%%%%%%%%%

\color{black}\noindent
        \item[\textbf{expected}:] assume that the \(k\) positions containing \(x\) are chosen uniformly at random
        among the \(n\) indices (or that the order of indices in the modular traversal is equally likely to ``hit'' any of these \(k\) with each step).
        Derive the \textbf{expected running time} in terms of \(n\) and \(k\) (asymptotically, e.g.\ \(\Theta(n/k)\) or similar),
        by introducing indicator random variables and using properties of expected value.

%%%%%%%%%%%%%%%%%%%%%%%%%%%%%%%%%%%%%%%%%%%%%%%%%
% Start of answer: Problem 3, Item 2(c)
%%%%%%%%%%%%%%%%%%%%%%%%%%%%%%%%%%%%%%%%%%%%%%%%%
\color{purple}
\begin{mdframed}
  \color{purple}
  \emph{Replace with your expected-time analysis.}
\end{mdframed}
%%%%%%%%%%%%%%%%%%%%%%%%%%%%%%%%%%%%%%%%%%%%%%%%%
% End of answer: Problem 3, Item 2(c)
%%%%%%%%%%%%%%%%%%%%%%%%%%%%%%%%%%%%%%%%%%%%%%%%%

      \end{enumerate}

\color{black}\noindent
    \item Consider a \emph{randomized} variant (\textsc{Randomized-Search}):
    instead of the fixed modular order, generate a uniformly random permutation \(\pi\) of \([1..n]\)
    and traverse the array by examining \(A[\pi(1)], A[\pi(2)], \dots, A[\pi(n)]\) until \(x\) is found.
      \begin{enumerate}[(a)]
        \item Write down pseudocode for \textsc{Randomized-Search} (including how you obtain or use the random permutation of indices).

%%%%%%%%%%%%%%%%%%%%%%%%%%%%%%%%%%%%%%%%%%%%%%%%%
% Start of answer: Problem 3, Item 3(a)
%%%%%%%%%%%%%%%%%%%%%%%%%%%%%%%%%%%%%%%%%%%%%%%%%
\color{purple}
\begin{mdframed}
  \color{purple}
  \emph{Replace with your pseudocode for \textsc{Randomized-Search}.}
\end{mdframed}
%%%%%%%%%%%%%%%%%%%%%%%%%%%%%%%%%%%%%%%%%%%%%%%%%
% End of answer: Problem 3, Item 3(a)
%%%%%%%%%%%%%%%%%%%%%%%%%%%%%%%%%%%%%%%%%%%%%%%%%

\color{black}\noindent
        \item Derive the \textbf{expected running time}, in terms of \(n\) and \(k\),
        for \textbf{any} fixed input array that contains exactly \(k\) copies of \(x\).

%%%%%%%%%%%%%%%%%%%%%%%%%%%%%%%%%%%%%%%%%%%%%%%%%
% Start of answer: Problem 3, Item 3(b)
%%%%%%%%%%%%%%%%%%%%%%%%%%%%%%%%%%%%%%%%%%%%%%%%%
\color{purple}
\begin{mdframed}
  \color{purple}
  \emph{Replace with your expected-time analysis for \textsc{Randomized-Search}.}
\end{mdframed}
%%%%%%%%%%%%%%%%%%%%%%%%%%%%%%%%%%%%%%%%%%%%%%%%%
% End of answer: Problem 3, Item 3(b)
%%%%%%%%%%%%%%%%%%%%%%%%%%%%%%%%%%%%%%%%%%%%%%%%%

      \end{enumerate}

\color{black}\noindent
    \item Compare \textsc{Modular-Search} vs.\ \textsc{Randomized-Search}: in which cases is the deterministic algorithm slow?
    How does randomizing the order make the algorithm more efficient in those cases? Argue in a short paragraph.

%%%%%%%%%%%%%%%%%%%%%%%%%%%%%%%%%%%%%%%%%%%%%%%%%
% Start of answer: Problem 3, Item 4
%%%%%%%%%%%%%%%%%%%%%%%%%%%%%%%%%%%%%%%%%%%%%%%%%
\color{purple}
\begin{mdframed}
  \color{purple}
  \emph{Replace with your comparison and discussion.}
\end{mdframed}
%%%%%%%%%%%%%%%%%%%%%%%%%%%%%%%%%%%%%%%%%%%%%%%%%
% End of answer: Problem 3, Item 4
%%%%%%%%%%%%%%%%%%%%%%%%%%%%%%%%%%%%%%%%%%%%%%%%%

  \end{enumerate}
\end{problem}

\begin{thebibliography}{CormenLeiserson09}

\bibitem[CLRS]{CLRS-2022}
Cormen, T.H., Leiserson, C.E., Rivest, R.L. and Stein, C., 2022. \emph{Introduction to Algorithms, Fourth Edition}. MIT Press.

\end{thebibliography}

\end{document}

